\documentclass[10pt,a4paper]{amsart}
\usepackage[margin=19mm]{geometry}
\usepackage{amsmath,amssymb,mathtools}
\usepackage[scr=rsfs,cal=euler]{mathalfa}
\usepackage{xfrac}
\usepackage[unicode,hidelinks,bookmarks=false]{hyperref}
\newcommand{\ie}{i.e.\ }
\newcommand{\cf}{cf.\ }
\newcommand{\resp}{respectively\ }
\newcommand{\Subsection}{\subsection}
\providecommand{\nobreakdash}{\nobreak\hbox{-}}
\setlength{\emergencystretch}{2em}
\multlinegap0pt
\usepackage{bm}
\usepackage{bbm}
\usepackage{dsfont}
\usepackage{xspace}
\usepackage{cancel}
\usepackage{mathtools}
\usepackage{amsthm}
\makeatletter
\@ifundefined{theorem}{\newtheorem{theorem}{Theorem}}{}
\@ifundefined{Fact}{\newtheorem{Fact}[theorem]{Fact}}{}
\makeatother
\def\bra#1{\left\langle#1\right|}
\def\gU{\textsf{U}}
\def\ket#1{\left|#1\right\rangle}
\title[Strict area law entanglement versus chirality]{Strict area law entanglement versus chirality}
\author{Xiang Li, Ting-Chun Lin, John McGreevy, Bowen Shi}
\date{Source: arXiv:2408.10306v3 (CC BY 4.0)}
\begin{document}
\maketitle

\noindent\emph{Source and licence.} Strict area law entanglement versus chirality. Version arXiv:2408.10306v3 (CC BY 4.0), distributed under the Creative Commons Attribution 4.0 International licence, as recorded in the arXiv licence metadata for this exact version and shown on the abstract page https://arxiv.org/abs/2408.10306v3. Full rights notice: licenses/arxiv-2408.10306-CC-BY-4.0.txt.\par\smallskip

\noindent\textit{Editorial scope.} Counted unit: the Fact whose source label is \texttt{fact:basic} in the article (source0.tex lines 1454--1462, source label \texttt{fact:basic}), delivered together with the article's own complete original proof of it (source0.tex lines 1463--1499, one \texttt{proof} environment of 37 lines). No mathematics is written, repaired or re-derived: statement and proof are the article's own text line for line. Required context retained uncounted: none. Every definition, display and label of those lines is retained unchanged. Omitted from this package and not redistributed: 1--1453, 1500--2031. Nothing inside a retained excerpt is omitted or altered: every retained body is delivered exactly as the article's lines are written, so no omission receipt of any kind accompanies this package. Citations: the retained excerpts carry no citation command at all, so no bibliography entry of the article is transcribed and no reference is redistributed. Reference adaptation: none was needed and none was made; every reference inside the delivered unit resolves inside it, so no unresolved reference or citation is produced. Printed slips kept literal: no printed word, symbol, hypothesis or number of the counted statement or of its proof was altered. Layout and class: the article's own class is \texttt{revtex4-2} with packages including geometry, titlesec, amssymb, amsmath, bm, graphicx, amsfonts, fancybox, xcolor, hyperref, soul, tikz, bclogo, pgffor; the wrapper is the lane's pinned \texttt{amsart} editorial wrapper at 10pt on A4, which restates the theorem-like environments the retained text needs and adds the article's own macro definitions that the retained text needs (\texttt{\textbackslash bra}, \texttt{\textbackslash gU}, \texttt{\textbackslash ket}), with extra wrapper packages (bm, bbm, dsfont, xspace, cancel, mathtools). The delivered numbering is the unit's own. Assets: no retained span contains a figure environment, a graphics-inclusion command, a PDF-page inclusion, a table or a TikZ picture; no image file is copied into this package, the package declares no asset dependency and no image file is redistributed with it. Licence: version arXiv:2408.10306v3 of this article is distributed under the Creative Commons Attribution 4.0 International licence, as recorded in the arXiv licence metadata for this exact version and shown on the abstract page https://arxiv.org/abs/2408.10306v3. A full rights notice is delivered at licenses/arxiv-2408.10306-CC-BY-4.0.txt. Characters: every retained excerpt is pure ASCII, so the delivered engine, XeLaTeX with its default Latin Modern text font, reports no missing character.
\begin{Fact}[Basic facts] \label{fact:basic}
    Below we collect four basic facts about the local charge operators of the on-site $\gU(1)$ symmetry and modular Hamiltonians from the symmetric state:  
    \begin{enumerate}
        \item $Q_{AB} = Q_A + Q_B$ for any regions $A,B$.  
        \item $[K_X,Q_X]=0$. Moreover, $[K_A,Q_{AB}] = 0$ and $ [K_{A},Q_{AB}^2] = 0$ for any regions $A,B$. 
        \item $\langle [K_{AB},Q_A^2]\rangle = \langle [K_{A},Q_{AB}^2]\rangle = 0$ for any region $AB$. 
        \item $\langle [K_{AB},Q_{BC}^2] \rangle = \langle [K_{AB},Q_{\overline{BC}}^2] \rangle$ for any tripartite region $ABC$. 
    \end{enumerate}
\end{Fact}

\begin{proof}
    Fact (1) is simply because the local charge operator is a sum over on-site operators. That is, for a region $X$, $Q_X = \sum_{v\in X} Q_v$. 

    Fact (2) is the consequence of the symmetry. Explicitly, $U(t)\ket{\Psi} = U_X(t)U_{\overline{X}}(t)\ket{\Psi} = \alpha \ket\Psi$, where $\alpha$ is a phase. Then $U_{X}(t)\ket{\Psi} = \alpha U_{\overline{X}}^{\dagger}(t)\ket{\Psi}$ and hence 
    \begin{equation}\label{eq:U-rho-commuting}\begin{split}
        &U_{X}(t)\ket{\Psi} \bra{\Psi}U_{X}(t)^{\dagger} = U_{\overline{X}}(t)^{\dagger}\ket{\Psi} \bra{\Psi}U_{\overline{X}}(t) \\ 
        \Rightarrow &\quad U_X(t)\rho_X U_X(t)^{\dagger} = \rho_X,
    \end{split}
    \end{equation}
    where the $\Rightarrow$ can be seen by tracing out $\overline{X}$. Then one can see $U_X(t)K_X U_X(t)^{\dagger} = K_X$ since $K_X = -\ln \rho_X$. Taking the first derivative of $t$ at $t=0$, we obtain $[K_X,Q_X]=0$. With $Q_{AB} = Q_{A}+Q_B$, we can see $[K_A,Q_{AB}]=0$. To see $[K_{A},Q_{AB}^2]= 0$, we can utilize the Fact (1) to write $Q_{AB} = Q_A + Q_B$. Therefore, 
    \begin{equation}\begin{split}
        & [K_{A},Q_{AB}^2]\\ 
        =& [K_{A},Q_{A}^2] +  [K_{A},Q_{B}^2] + 2  [K_{A},Q_{A}Q_B] \\ 
        =&0.
    \end{split}
    \end{equation}
    The three commutators above all vanish because $[K_A,Q_A]=0$ and $[K_A,Q_B]=0$.  
        
    Now we explain Fact (3). To see $\langle [K_{AB},Q_A^2]\rangle = 0$, we utilize the flipping $K_{AB}\ket\Psi = K_{\overline{AB}}\ket{\Psi}$ due to the property of pure state. Therefore, 
    \begin{equation}
        \langle [K_{AB},Q_A^2]\rangle = \langle [K_{\overline{AB}},Q_A^2]\rangle = 0,
    \end{equation}
    as the compliment of $AB$ has no overlap with $A$. 

    Fact (4) is for any tripartite region $ABC$. Because the state has $\gU(1)$ symmetry, $\ket\Psi$ is an eigenstate of the symmetry operator $Q$. Let $q$ be the eigenvalue, we thus have $Q\ket{\Psi} = (Q_{BC}+Q_{\overline{BC}})\ket{\Psi} = q \ket\Psi$. Therefore, we can replace $Q_{BC}\ket\Psi$ with $(q-Q_{\overline{BC}})\ket\Psi$ and as a result 
    \begin{equation}\begin{split}
        \langle [K_{AB},Q_{BC}^2] \rangle &= \langle [K_{AB},q^2-2q Q_{\overline{BC}} + Q_{\overline{BC}}^2]\rangle \\ 
        &= \langle [K_{AB}, Q_{\overline{BC}}^2]\rangle. 
    \end{split}
    \end{equation}
    In the derivation of the last line in the equation above, $\langle [K_{AB},q^2]\rangle = 0$ is obvious as $q$ is simply a number. To see $\langle [K_{AB},2q Q_{BC}]\rangle = 0$, we can show that 
    \begin{equation}\begin{split}
        \langle [K_{AB},Q_{BC}]\rangle =& \langle [K_{AB},Q_{B}]\rangle + \langle [K_{AB},Q_{C}]\rangle  \\ 
        =& \langle [K_{\overline{AB}},Q_{B}]\rangle = 0. 
    \end{split}
    \end{equation} 
\end{proof}

\end{document}
